%% kapa: exam=1 round=32 year=2021 period=1 form=A subject=real_estate subject_ko=부동산학원론 seq=35 number=115
%% source: Q-Net 감정평가사 기출문제 (cst00309) artlSeq=5208397 file=extracted/1008/5208397/A형/2021년도 제32회 감정평가사 1차 1교시 A형.pdf
%% conversion: kordoc 4.24.51 -> tools/gen1.py (auto); review=proofread-sonnet; flags=none
\begin{question}{115}
부동산 평가활동에서 부동산 가격의 원칙에 관한 설명으로 옳지 않은 것은?
\begin{choices}
  \item 예측의 원칙이란 평가활동에서 가치형성요인의 변동추이 또는 동향을 주시해야 한다는 것을 말한다.
  \item 대체의 원칙이란 부동산의 가격이 대체관계의 유사 부동산으로부터 영향을 받는다는 것을 말한다.
  \item 균형의 원칙이란 부동산의 유용성이 최고도로 발휘되기 위해서는 부동산이 외부환경과 균형을 이루어야 한다는 것을 말한다.
  \item 변동의 원칙이란 가치형성요인이 시간의 흐름에 따라 지속적으로 변화함으로써 부동산 가격도 변화한다는 것을 말한다.
  \item 기여의 원칙이란 부동산의 가격이 대상부동산의 각 구성요소가 기여하는 정도의 합으로 결정된다는 것을 말한다.
\end{choices}
\end{question}
\answer{3}
